%%%%%%%%%%%%%%%%%
% This is the resume editor's preamble for rt-altacv.cls, a modified version of the sample
% CV altacv-sample.tex (v1.7.4, 30 July 2025) that LianTze Lim (liantze@gmail.com) wrote
% to go with altacv.cls. It is not the original, which is at
% https://github.com/liantze/AltaCV and on CTAN as altacv.
%
%% It may be distributed and/or modified under the
%% conditions of the LaTeX Project Public License, either version 1.3
%% of this license or (at your option) any later version.
%% The latest version of this license is in
%%    http://www.latex-project.org/lppl.txt
%% and version 1.3 or later is part of all distributions of LaTeX
%% version 2003/12/01 or later.
%%%%%%%%%%%%%%%%

% This is the editor's port of AltaCV. The layout, the fonts pdfLaTeX takes, the colours and the
% fonts of the name and the section titles are the sample's, and the class is AltaCV's with only
% what the editor's compile loop needs changed, as rt-altacv.cls says at its top. The port changes
% the sample only as follows, each where it is made. The class is the renamed one, and takes no
% withhyper option, since it always loads hyperref. No font sets ligatures, and the name and the
% section titles are not set in capitals, because the editor's page cannot yet match either to the
% letters the resume holds. A symbol of LaTeX's TS1 encoding that a family lacks, such as an arrow
% in Roboto Slab, is set in Latin Modern, where LaTeX would take Computer Modern, whose TS1 fonts
% the editor's TeX Live has no outlines for, so that neither the page nor the PDF could show it. The
% margins control sets the margins after the format, the same on every side, and its default is half
% an inch, the nearest it has to the sample's 1.25 cm at the sides and 1.5 cm at the top and bottom,
% so only the space between the columns is set here. The accent colour is named, rtaccent, and the
% sample's red is its default; the colours the sample sets from that red, and the section titles'
% and their rules', are set from rtaccent by \rtaccentcolours, so that the accent control recolours
% the tagline, the icons, the entries' subtitles and the section titles with their rules. The
% sample's rule is a gold no accent can derive, so it becomes a tint of the accent. The class sets
% its sizes once, as it loads, from LaTeX's size10.clo, while the font size control takes half
% points from 9 to 12 as an ordinary compile: \rtfontsize sets those sizes and the spacing that goes
% with them scaled by the size asked for over 10 points, and \rtdensity scales the space between
% sections, entries and bullets. At 10 points and normal density, the controls' defaults, every
% length is the class's. The photo, the publications and the wheel chart are left out, since the
% editor's model has no place for them. And the file ends where the body would begin, because the
% editor generates the body from the resume.

\documentclass[10pt,a4paper]{rt-altacv}

% The sample's space between the columns. The margins control sets the margins after the format.
\geometry{columnsep=1.2cm}

% The paracol package lets you typeset columns of text in parallel
\usepackage{paracol}

% The fonts the sample takes on pdfLaTeX: Roboto Slab as the roman family, which the name and the
% section titles are set in, and Lato as the sans-serif family, which is the text's. The font
% control sets the text's family again after the format.
\usepackage[rm]{roboto}
\usepackage[defaultsans]{lato}
\renewcommand{\familydefault}{\sfdefault}
% TS1 symbols a family lacks, such as arrows, come from Latin Modern, which has their outlines.
\def\rmsubstdefault{lmr}
\def\sfsubstdefault{lmss}
\def\ttsubstdefault{lmtt}
% No ligatures, in any font. Lato names its ff, fi, fl, ffi and ffl glyphs uniFB00 to uniFB04,
% Roboto the last four uniFB01 to uniFB04, and Roboto Slab its fl, ffi and ffl, which the editor's
% page reads as the single characters of Unicode's presentation forms, so that a field with fi in
% it would not align with the letters the resume holds, and a click in it would land on the wrong
% one. So each font is set without its ligatures as it is selected, and a user's -- stays two
% hyphens. This can go once the page reads those names as the letters they join, as it reads the
% ones Latin Modern names fi.
\AddToHook{selectfont}{\pdfnoligatures\font}

% The font size control's sizes, and the density control's spacing, as the note at the top says.
% \rt@size is the size in scaled points, 10 points until \rtfontsize sets another, and
% \rt@pt{<points>} is a length of size10.clo's scaled by it over 10 points and rounded to a
% hundredth of a point, TeX arithmetic throughout. A length to two decimals rounds back to what TeX
% reads for it, so at 10 points every length is the class's.
\makeatletter
\newcommand*\rt@size{655360}
\newcommand*\rt@pt[1]{%
  \the\dimexpr\numexpr\dimexpr#1\p@*\rt@size/655360\relax*100/65536\relax\p@/100\relax}
% size10.clo's sizes, with the display skips and the lists' spacing it sets with them, each length
% through \rt@pt.
\DeclareRobustCommand\normalsize{%
  \@setfontsize\normalsize{\rt@pt{10}}{\rt@pt{12}}%
  \abovedisplayskip \rt@pt{10}\@plus\rt@pt{2}\@minus\rt@pt{5}%
  \abovedisplayshortskip \z@ \@plus\rt@pt{3}%
  \belowdisplayshortskip \rt@pt{6}\@plus\rt@pt{3}\@minus\rt@pt{3}%
  \belowdisplayskip \abovedisplayskip
  \let\@listi\@listI}
\DeclareRobustCommand\small{%
  \@setfontsize\small{\rt@pt{9}}{\rt@pt{11}}%
  \abovedisplayskip \rt@pt{8.5}\@plus\rt@pt{3}\@minus\rt@pt{4}%
  \abovedisplayshortskip \z@ \@plus\rt@pt{2}%
  \belowdisplayshortskip \rt@pt{4}\@plus\rt@pt{2}\@minus\rt@pt{2}%
  \def\@listi{\leftmargin\leftmargini
    \topsep \rt@pt{4}\@plus\rt@pt{2}\@minus\rt@pt{2}%
    \parsep \rt@pt{2}\@plus\rt@pt{1}\@minus\rt@pt{1}%
    \itemsep \parsep}%
  \belowdisplayskip \abovedisplayskip}
\DeclareRobustCommand\footnotesize{%
  \@setfontsize\footnotesize{\rt@pt{8}}{\rt@pt{9.5}}%
  \abovedisplayskip \rt@pt{6}\@plus\rt@pt{2}\@minus\rt@pt{4}%
  \abovedisplayshortskip \z@ \@plus\rt@pt{1}%
  \belowdisplayshortskip \rt@pt{3}\@plus\rt@pt{1}\@minus\rt@pt{2}%
  \def\@listi{\leftmargin\leftmargini
    \topsep \rt@pt{3}\@plus\rt@pt{1}\@minus\rt@pt{1}%
    \parsep \rt@pt{2}\@plus\rt@pt{1}\@minus\rt@pt{1}%
    \itemsep \parsep}%
  \belowdisplayskip \abovedisplayskip}
\DeclareRobustCommand\scriptsize{\@setfontsize\scriptsize{\rt@pt{7}}{\rt@pt{8}}}
\DeclareRobustCommand\tiny{\@setfontsize\tiny{\rt@pt{5}}{\rt@pt{6}}}
\DeclareRobustCommand\large{\@setfontsize\large{\rt@pt{12}}{\rt@pt{14}}}
\DeclareRobustCommand\Large{\@setfontsize\Large{\rt@pt{14.4}}{\rt@pt{18}}}
\DeclareRobustCommand\LARGE{\@setfontsize\LARGE{\rt@pt{17.28}}{\rt@pt{22}}}
\DeclareRobustCommand\huge{\@setfontsize\huge{\rt@pt{20.74}}{\rt@pt{25}}}
\DeclareRobustCommand\Huge{\@setfontsize\Huge{\rt@pt{24.88}}{\rt@pt{30}}}
% size10.clo's list spacing: the first level's at the normal size, which \normalsize hands to
% \@listi, and the second and third levels'. The class's lists take enumitem's spacing over them.
\def\rt@listI{%
  \parsep \rt@pt{4}\@plus\rt@pt{2}\@minus\rt@pt{1}%
  \topsep \rt@pt{8}\@plus\rt@pt{2}\@minus\rt@pt{4}%
  \itemsep \rt@pt{4}\@plus\rt@pt{2}\@minus\rt@pt{1}}
\def\@listI{\leftmargin\leftmargini\rt@listI}
\def\@listii{\leftmargin\leftmarginii
  \labelwidth\leftmarginii
  \advance\labelwidth-\labelsep
  \topsep \rt@pt{4}\@plus\rt@pt{2}\@minus\rt@pt{1}%
  \parsep \rt@pt{2}\@plus\rt@pt{1}\@minus\rt@pt{1}%
  \itemsep \parsep}
\def\@listiii{\leftmargin\leftmarginiii
  \labelwidth\leftmarginiii
  \advance\labelwidth-\labelsep
  \topsep \rt@pt{2}\@plus\rt@pt{1}\@minus\rt@pt{1}%
  \parsep \z@
  \partopsep \rt@pt{1}\@plus\z@ \@minus\rt@pt{1}%
  \itemsep \topsep}
% The space the class's macros leave is size10.clo's \smallskip, \medskip and \bigskip: \cvsection
% leaves a \bigskip before a title and a \medskip after its rule, \cvevent a \smallskip after its
% title and its subtitle and a \medskip after its dates, \divider a \medskip, and \makecvheader a
% \medskip after each of its lines. \rt@density is a percentage of them, 100 until \rtdensity sets
% another, and \rt@gap{<points>} is a length through \rt@pt, scaled by it and rounded again.
\newcommand*\rt@density{100}
\newcommand*\rt@gap[1]{%
  \the\dimexpr\numexpr\dimexpr\rt@pt{#1}*\rt@density/100\relax*100/65536\relax\p@/100\relax}
\newcommand*\rt@skips{%
  \setlength\smallskipamount{\rt@gap{3}\@plus\rt@gap{1}\@minus\rt@gap{1}}%
  \setlength\medskipamount{\rt@gap{6}\@plus\rt@gap{2}\@minus\rt@gap{2}}%
  \setlength\bigskipamount{\rt@gap{12}\@plus\rt@gap{4}\@minus\rt@gap{4}}}
% \rtfontsize{<points>}: every size and its spacing from <points> in place of 10, then the normal
% size, the first list level's spacing and \partopsep outside any list, where \trivlist takes them,
% and the skips, as size10.clo sets them as it loads. The paragraph indent stays the class's 0pt.
\newcommand*\rtfontsize[1]{%
  \edef\rt@size{\number\dimexpr#1\p@\relax}%
  \normalsize
  \rt@listI
  \partopsep \rt@pt{2}\@plus\rt@pt{1}\@minus\rt@pt{1}\relax
  \rt@skips}
% \rtdensity{<percent>}: the skips at that share of what the size gives them, and the space
% enumitem leaves between items and after a list, which the class sets to a quarter of a line, at
% that share of a quarter.
\newcommand*\rtdensity[1]{%
  \edef\rt@density{\number#1}%
  \rt@skips
  \edef\rt@items{\strip@pt\dimexpr0.25\p@*\rt@density/100\relax}%
  \setlist*{itemsep=\rt@items\baselineskip,after=\vspace{\rt@items\baselineskip}}}
\makeatother

% The sample's colours. The accent is its red, named rtaccent; the accent control defines it again
% after the format and sets \rtaccentcolours again, since \colorlet takes a colour's value when it
% runs. The name, the entries' titles and the text keep the sample's colours.
\definecolor{SlateGrey}{HTML}{2E2E2E}
\definecolor{LightGrey}{HTML}{666666}
\definecolor{rtaccent}{HTML}{8F0D0D}
\colorlet{name}{black}
\colorlet{emphasis}{SlateGrey}
\colorlet{body}{LightGrey}
% The colours the sample sets from its red, from the accent: the tagline, the icons and the
% entries' subtitles in it, the section titles in it darkened by half, which is the sample's
% DarkPastelRed to within a few units, and the rules under them in a light tint of it.
\newcommand*\rtaccentcolours{%
  \colorlet{accent}{rtaccent}%
  \colorlet{tagline}{rtaccent}%
  \colorlet{subheading}{rtaccent}%
  \colorlet{heading}{rtaccent!50!black}%
  \colorlet{headingrule}{rtaccent!30}}
\rtaccentcolours

% Change some fonts, if necessary
\renewcommand{\namefont}{\Huge\rmfamily\bfseries}
\renewcommand{\personalinfofont}{\footnotesize}
\renewcommand{\cvsectionfont}{\LARGE\rmfamily\bfseries}
\renewcommand{\cvsubsectionfont}{\large\bfseries}

% The name and the section titles as the user types them. The class sets both in capitals through
% \MakeUppercase, which the port takes out of \makecvheader and \cvsection: the editor's page
% cannot yet match a capital it draws to the lower-case letter the resume holds, so a click would
% land on the wrong one; and \MakeUppercase takes time that grows with the square of its argument's
% length, about 2 seconds a pass for a title of 8,000 characters. In \cvsection, whose optional
% argument makes \cvsection a wrapper around the macro that sets the title, \leavevmode takes its
% place, so that a blank title still starts the line the class ends with \\, which would otherwise
% stop the compile. A class that no longer has them stops the dump.
\makeatletter
\newcommand*\rt@unpatched[1]{\errmessage{rt-altacv: no \string\MakeUppercase\space in \string#1}}
\patchcmd{\makecvheader}{\MakeUppercase}{}{}{\rt@unpatched\makecvheader}
\expandafter\patchcmd\csname\string\cvsection\endcsname{\MakeUppercase}{\leavevmode}{}%
  {\rt@unpatched\cvsection}
\makeatother
